\documentclass[11pt,a4paper]{article}
\usepackage{iftex}
\ifPDFTeX
  \usepackage[T1]{fontenc}
  \usepackage[utf8]{inputenc}
  \usepackage{lmodern}
  \input{glyphtounicode}
  \pdfgentounicode=1
\else
  \usepackage{fontspec}
  % Kerning off: kerned pairs ("A WS", "T ree") split words for ATS text extraction.
  \setmainfont{lmroman10-regular}[Extension=.otf, BoldFont=lmroman10-bold,
    ItalicFont=lmroman10-italic, BoldItalicFont=lmroman10-bolditalic, Kerning=Off]
\fi
\usepackage[margin=0.9in]{geometry}
\usepackage[hidelinks]{hyperref}
\hypersetup{pdftitle={Abhinav Kaushik - Cover Letter - Signifyd}, pdfauthor={Abhinav Kaushik}}

\pagestyle{empty}
\setlength{\parindent}{0pt}
\setlength{\parskip}{10pt}
\raggedright

\begin{document}

{\Large\bfseries Abhinav Kaushik}\\
Data Scientist II\\
New Delhi, India \textbar{} +91 8700571859 \textbar{} \href{mailto:aabhinav.kaushik@gmail.com}{aabhinav.kaushik@gmail.com} \textbar{} \href{https://www.linkedin.com/in/abhinav-kaushik10/}{LinkedIn}  

September 27, 2026

Hiring Team\\
Signifyd

\textbf{Re: Data Scientist II}

Dear Hiring Team,

I am applying for the Data Scientist II role in the Applied Decision Science team. Most of my work has been building prediction models and taking them all the way into production, from SLA breach prediction in telecom to analytics tools in healthcare. I like owning the full loop: the data, the model, the deployment and the numbers afterwards.

Fraud decisioning is a problem where a model's choices touch real customers every day, and a team that runs its own experiments end to end is where I do my best work. It also fits my life: my fiancée lives in Berlin, and we plan to move to London together and settle there long term.

I would bring a habit of shipping models, not just notebooks. I deployed a LightGBM breach prediction system with FastAPI and MLflow that reduced downtime by 67\%, and earlier built an XGBoost propensity model that helped a lending team prioritize customers. I am comfortable in Python, SQL and PySpark on Databricks.

Beyond my day job I am building ThinkTree.AI, a non-profit AI learning platform used by NGO teachers to support students with ADHD. Working with teachers taught me to test ideas against what actually helps people. I would welcome a conversation about the role.

Sincerely,\\[4pt]
Abhinav Kaushik

\end{document}